\color{unidarkgreen}{\subsection{Управление выключателем}}
\color{black}

\begin{enumerate}[label=\arabic{section}.\arabic{subsection}.\arabic*,  align=left, labelsep=4pt, leftmargin=0pt, itemindent=57pt]

\item
Функциональный блок управления выключателем (УВ) предназначен для подачи сигнала на включение
силового выключателя. 


\begin{figure}[H]
\centering
\includegraphics[width=0.8\textwidth,height=0.8\textheight,keepaspectratio]{\fbpath/761.3431 УВ/img/UV.pdf}
\captionof{figure}{Функциональная схема УВ}\label{pic_uv:UV}
\end{figure}

\item
Команда на включение выключателя \textit{«Включить В»} формируется:
\begin{itemize}
\item от команды оперативного включения выключателя (сигнал \textit{«РКВ»});
\item при срабатывании АПВ (сигнал \textit{«Ср. АПВ»}).
\end{itemize}

\item
Ввод в работу логики управления выключателем осуществляется с помощью программного ключа \textbf{«УВ»}.
Оперативное включение выключателя осуществляется с контролем наличия разрешающего сигнала от
функционального блока КСН (принцип работы которого описан {\color{red}в пункте 13} настоящего руководства по
эксплуатации). Максимальное время, в течение которого продолжается ожидание условий оперативного
включения, регулируется уставкой \textbf{«Tожид усл опер вкл»}. По истечению данной выдержки времени команда
оперативного включения сбрасывается. Выдержка времени на возврат \textbf{«Тпрод опер вкл»} задает время
продления сигнала оперативного включения.

\item
При необходимости может быть введен подхват команды на включение от токового реле,
установленного в цепи ЭМВ, в этом случае сигнал на включение удерживается в сработанном состоянии до
исчезновения тока через ЭМВ (пока блок-контакты выключателя не разорвут цепь включения). Ввод
осуществляется программным ключом \textbf{«Подхват вкл В от РТ ЭМВ»}. 

\item
Предусмотрена возможность блокировки многократных включений на КЗ. Если при наличии сигнала
включения будет протекать ток по электромагнитам отключения (что соответствует включению на КЗ), то
сигнал \textit{«Включить В»} заблокируется, пока не пропадут сигналы включения (от АПВ или оперативного
включения). Блокировка вводится программным ключом \textbf{«Блок от многокр вкл на КЗ»}.

\item
При наличии сигнала \textit{«Блок. включения В»}, формируемого функциональным блоком контроля силового
выключателя, команда на включение выключателя блокируется.

\item
Оперативный ввод/вывод УВ производится с помощью БУ (блоков управления функциями) \textit{«Вывод УВ»}. 

\small
\begin{longtable}{|>{\centering\arraybackslash}m{5.3cm}|>{\centering\arraybackslash}m{4.2cm}|>{\centering\arraybackslash}m{1cm}|>{\centering\arraybackslash}m{3.3cm}|>{\centering\arraybackslash}m{1.8cm}|}
% Первая страница: полный заголовок
\caption{Перечень уставочных параметров УВ\hfill\vspace{-0.5\baselineskip}}\label{uv:tbl1}\\
\hline
\rowcolor{gray!30}
Наименование уставки & Диапазон & Шаг & Значение по умолчанию & Единица измерения \\
\hline
\endfirsthead

% Промежуточные страницы: «Продолжение таблицы»
\caption*{\hspace{3pt}\emph{Продолжение таблицы \ref{uv:tbl1}}\hfill\ }\vspace{-0.5\baselineskip} \\
\hline
\rowcolor{gray!30}
Наименование уставки & Диапазон & Шаг & Значение по умолчанию & Единица измерения \\
\hline
\endhead

%===m>УВ|-
%===m
\end{longtable}
\normalsize	

\end{enumerate}